\documentclass[aspectratio=169]{beamer}

\usetheme{Madrid}
\usecolortheme{default}
\setbeamertemplate{navigation symbols}{}
\setbeamertemplate{footline}[frame number]
\setbeamersize{text margin left=0.45cm, text margin right=0.45cm}

\usepackage[T1]{fontenc}
\usepackage[utf8]{inputenc}
\usepackage{graphicx}
\usepackage{booktabs}
\usepackage{xcolor}
\usepackage{tikz}
\usepackage{ifthen}

% ============================================================
% User configuration
% ============================================================
\newcommand{\Campaign}{2022F}
\newcommand{\SelectionMode}{symmetric}
\newcommand{\PeriodLabel}{Run \Campaign}
\newcommand{\DataTightDir}{DATA_tight}
\newcommand{\DataRelaxedDir}{DATA_relaxed}
\newcommand{\TightLabel}{tight}
\newcommand{\RelaxedLabel}{relaxed}

% Edit these folder paths to point to the analysis outputs you want to present.
\newcommand{\DGLBaseDir}{../results_full_scan_DGL_2022F_symmetric}
\newcommand{\DSABaseDir}{../results_full_scan_DSA_2022F_symmetric}
\newcommand{\DGLBaseDirA}{../results_full_scan_DGL_2022F_symmetric}
\newcommand{\DGLBaseDirB}{../results_full_scan_DGL_2023D_symmetric}
\newcommand{\DGLBaseDirC}{../results_full_scan_DGL_2024I_symmetric}

% ============================================================
% Helpers
% ============================================================
\newcommand{\tightLooseTable}{
  \begin{tabular}{@{}lll@{}}
    \toprule
    Quantity & Relaxed & Tight \\
    \midrule
    Tag/probe pairing & Upstream ntuplizer & Same upstream pairing \\
    $p_{T}^{err}/p_{T}$ (DGL) & not applied & $< 0.3$ on tag and probe \\
    $\chi^{2}$ (DGL) & not applied & $< 5$ on tag and probe \\
    MuonHits (DGL) & not applied & $\geq 13$ on tag and probe \\
    ValidStripHits (DGL) & not applied & $\geq 6$ on tag and probe \\
    $p_{T}^{err}/p_{T}$ (DSA) & not applied & $< 0.2$ on tag and probe \\
    $\chi^{2}$ (DSA) & not applied & $< 5$ on tag and probe \\
    ValidMuonDTHits (DSA) & not applied & $\geq 31$ on tag and probe \\
    \bottomrule
  \end{tabular}
}

\newcommand{\legacyTable}{
  \begin{tabular}{@{}lll@{}}
    \toprule
    Scenario & Tag & Probe \\
    \midrule
    Baseline & no extra cut & no extra cut \\
    TagOnly & $p_{T}^{err}/p_{T} < 0.2$ & no extra cut \\
    TagProbe & $p_{T}^{err}/p_{T} < 0.2$ & $p_{T}^{err}/p_{T} < 0.2$ \\
    \bottomrule
  \end{tabular}
}

\newcommand{\plotornote}[3][0.60\textheight]{%
  \IfFileExists{#2}{%
    \includegraphics[width=\linewidth,height=#1,keepaspectratio]{#2}%
  }{%
    \begin{tikzpicture}
      \draw[rounded corners=4pt, thick, gray!70] (0,0) rectangle (11.6,5.8);
      \node[align=center, text width=10.4cm] at (5.8,2.9) {\Large Missing plot\\[0.4em]\normalsize #3\\[0.6em]\scriptsize\ttfamily\detokenize{#2}};
    \end{tikzpicture}%
  }%
}

\newcommand{\twoplots}[4]{%
  \begin{columns}[T,onlytextwidth]
    \column{0.49\textwidth}
    \centering
    \plotornote[0.52\textheight]{#1}{#2}
    \column{0.49\textwidth}
    \centering
    \plotornote[0.52\textheight]{#3}{#4}
  \end{columns}
}

\newcommand{\singleplot}[2]{%
  \begin{center}
    \plotornote[0.60\textheight]{#1}{#2}
  \end{center}
}

\newcommand{\twoplotsplusone}[6]{%
  \twoplots{#1}{#2}{#3}{#4}
  \vspace{0.3em}
  \begin{center}
    \plotornote[0.22\textheight]{#5}{#6}
  \end{center}
}

\newcommand{\threebulletsummary}[3]{%
  \begin{itemize}
    \item #1
    \item #2
    \item #3
  \end{itemize}
}

\title{Cosmic Muon Resolution Studies}
\subtitle{DGL Results by Year and DSA Bias Status}
\author{Alfredo Casta\~neda}
\date{July 19, 2026}

\begin{document}

\begin{frame}
  \titlepage
\end{frame}

\begin{frame}{Agenda}
  \tableofcontents
\end{frame}

\section{Selection Strategy}

\begin{frame}{Goal of This Update}
  \threebulletsummary
    {Establish the DGL result as the main physics message, consistently across the 2022F, 2023D, and 2024I campaigns.}
    {Use the symmetric workflow and control plots to show that the DGL selection is stable and well behaved.}
    {Report the current DSA bias status separately, together with the upstream ntuplizer study plan.}
\end{frame}

\begin{frame}{Scenario Definitions}
  \begin{block}{Symmetric workflow}
    \tightLooseTable
  \end{block}
  \vspace{0.5em}
  \begin{itemize}
    \item The tag/probe identity still comes from the ntuplizer.
    \item The tight scenario adds the same quality cuts to both legs of the pair.
    \item A single symmetric production run writes both \texttt{DATA\_relaxed}/\texttt{MC\_relaxed} and \texttt{DATA\_tight}/\texttt{MC\_tight}.
  \end{itemize}
\end{frame}

\begin{frame}{Legacy Reference}
  \begin{block}{Original three-scenario study}
    \legacyTable
  \end{block}
  \vspace{0.5em}
  \begin{itemize}
    \item Useful as a historical baseline.
    \item The symmetric workflow is intended to reduce tag/probe quality bias.
  \end{itemize}
\end{frame}

\section{DGL Controls}

\begin{frame}{What To Check in The Control Plots}
  \threebulletsummary
    {If tag and probe show different $p_{T}^{err}/p_{T}$ spectra, first check whether their $p_{T}$ spectra are already different.}
    {The profile plot of $p_{T}^{err}/p_{T}$ versus $p_{T}$ helps separate a kinematic effect from a pure selection effect.}
    {The hit-count and $\chi^{2}$ overlays test whether the tight scenario really balances the track-quality distributions.}
\end{frame}

\begin{frame}{DGL: Tag vs Probe $p_{T}$ and $p_{T}^{err}/p_{T}$}
  \twoplots
    {\DGLBaseDir/\DataTightDir/DATA/DGL/Control_plots/pt_tag_probe_comparison.png}
    {DGL DATA tight: tag vs probe $p_{T}$}
    {\DGLBaseDir/\DataTightDir/DATA/DGL/Control_plots/ptErr_over_pt_comparison.png}
    {DGL DATA tight: tag vs probe $p_{T}^{err}/p_{T}$}
\end{frame}

\begin{frame}{DGL: Residual Distribution Comparison}
  \twoplots
    {\DGLBaseDir/\DataTightDir/DATA/DGL/Control_plots/Tot_pthist.png}
    {DGL DATA tight: standard residual}
    {\DGLBaseDir/\DataTightDir/DATA/DGL/Control_plots/Tot_pthist_sym.png}
    {DGL DATA tight: symmetric residual}
\end{frame}

\begin{frame}{DGL: Residual Distributions vs $p_{T}$}
  \twoplots
    {\DGLBaseDir/\DataTightDir/DATA/DGL/pt_Plots/DGL_pt_all_fits.png}
    {DGL DATA tight: standard residual grid vs $p_{T}$}
    {\DGLBaseDir/\DataTightDir/DATA/DGL/pt_Plots/DGL_pt_all_fits_sym.png}
    {DGL DATA tight: symmetric residual grid vs $p_{T}$}
\end{frame}

\begin{frame}{DGL: Is $p_{T}^{err}/p_{T}$ Difference Mainly Kinematic?}
  \twoplots
    {\DGLBaseDir/\DataTightDir/DATA/DGL/Control_plots/tag_ptErrOverPt_vs_pt.png}
    {DGL DATA tight: tag $p_{T}^{err}/p_{T}$ vs $p_{T}$}
    {\DGLBaseDir/\DataTightDir/DATA/DGL/Control_plots/probe_ptErrOverPt_vs_pt.png}
    {DGL DATA tight: probe $p_{T}^{err}/p_{T}$ vs $p_{T}$}
\end{frame}

\begin{frame}{DGL: $p_{T}^{err}/p_{T}$ Profile Comparison}
  \singleplot
    {\DGLBaseDir/\DataTightDir/DATA/DGL/Control_plots/ptErrOverPt_vs_pt_profile_comparison.png}
    {DGL DATA tight: profile comparison}
\end{frame}

\begin{frame}{DGL: Track-Quality Control Plots}
  \twoplots
    {\DGLBaseDir/\DataTightDir/DATA/DGL/Control_plots/chi2_comparison.png}
    {DGL DATA tight: $\chi^{2}$ comparison}
    {\DGLBaseDir/\DataTightDir/DATA/DGL/Control_plots/MuonHits_comparison.png}
    {DGL DATA tight: MuonHits comparison}
\end{frame}

\begin{frame}{DGL: ValidStripHits Control Plot}
  \singleplot
    {\DGLBaseDir/\DataTightDir/DATA/DGL/Control_plots/ValidStripHits_comparison.png}
    {DGL DATA tight: ValidStripHits comparison}
\end{frame}

\section{DGL Results}

\begin{frame}{Reference Category}
  \threebulletsummary
    {The DGL headline plots are based on the tight symmetric category.}
    {This category applies the same downstream quality cuts to both tag and probe and is our default for year-to-year reporting.}
    {Relaxed results remain available as cross-check material, but the main DGL message should be built from one consistent category across all campaigns.}
\end{frame}

\begin{frame}{Symmetric Residual Diagnostic}
  \begin{block}{Definition}
    \[
      R_{\mathrm{sym}} =
      \frac{(q/p_{T})_{\mathrm{probe}} - (q/p_{T})_{\mathrm{tag}}}
      {\tfrac{1}{2}\left[(q/p_{T})_{\mathrm{probe}} + (q/p_{T})_{\mathrm{tag}}\right]}
    \]
  \end{block}
  \vspace{0.4em}
  \begin{itemize}
    \item Unlike the standard residual, the denominator does not privilege the tag leg alone.
    \item This is a diagnostic observable used to test whether the bias comes from the residual definition or from a real tag/probe reconstruction asymmetry.
    \item If the standard and symmetric residuals are both shifted, the asymmetry is likely intrinsic to the reconstructed pair.
  \end{itemize}
\end{frame}

\begin{frame}{DGL 2022F: MC vs Data Mean and Sigma}
  \twoplots
    {\DGLBaseDirA/compare_MC_DATA_DGL_\TightLabel/pt_mean_mc_vs_data.png}
    {DGL 2022F tight: $p_{T}$ mean, MC vs Data}
    {\DGLBaseDirA/compare_MC_DATA_DGL_\TightLabel/pt_sigma_mc_vs_data.png}
    {DGL 2022F tight: $p_{T}$ sigma, MC vs Data}
\end{frame}

\begin{frame}{DGL 2022F: Symmetric Residual Diagnostic}
  \twoplots
    {\DGLBaseDirA/compare_MC_DATA_DGL_\TightLabel/pt_mean_sym_mc_vs_data.png}
    {DGL 2022F tight: symmetric residual mean}
    {\DGLBaseDirA/compare_MC_DATA_DGL_\TightLabel/pt_sigma_sym_mc_vs_data.png}
    {DGL 2022F tight: symmetric residual sigma}
\end{frame}

\begin{frame}{DGL 2023D: MC vs Data Mean and Sigma}
  \twoplots
    {\DGLBaseDirB/compare_MC_DATA_DGL_\TightLabel/pt_mean_mc_vs_data.png}
    {DGL 2023D tight: $p_{T}$ mean, MC vs Data}
    {\DGLBaseDirB/compare_MC_DATA_DGL_\TightLabel/pt_sigma_mc_vs_data.png}
    {DGL 2023D tight: $p_{T}$ sigma, MC vs Data}
\end{frame}

\begin{frame}{DGL 2023D: Symmetric Residual Diagnostic}
  \twoplots
    {\DGLBaseDirB/compare_MC_DATA_DGL_\TightLabel/pt_mean_sym_mc_vs_data.png}
    {DGL 2023D tight: symmetric residual mean}
    {\DGLBaseDirB/compare_MC_DATA_DGL_\TightLabel/pt_sigma_sym_mc_vs_data.png}
    {DGL 2023D tight: symmetric residual sigma}
\end{frame}

\begin{frame}{DGL 2024I: MC vs Data Mean and Sigma}
  \twoplots
    {\DGLBaseDirC/compare_MC_DATA_DGL_\TightLabel/pt_mean_mc_vs_data.png}
    {DGL 2024I tight: $p_{T}$ mean, MC vs Data}
    {\DGLBaseDirC/compare_MC_DATA_DGL_\TightLabel/pt_sigma_mc_vs_data.png}
    {DGL 2024I tight: $p_{T}$ sigma, MC vs Data}
\end{frame}

\begin{frame}{DGL 2024I: Symmetric Residual Diagnostic}
  \twoplots
    {\DGLBaseDirC/compare_MC_DATA_DGL_\TightLabel/pt_mean_sym_mc_vs_data.png}
    {DGL 2024I tight: symmetric residual mean}
    {\DGLBaseDirC/compare_MC_DATA_DGL_\TightLabel/pt_sigma_sym_mc_vs_data.png}
    {DGL 2024I tight: symmetric residual sigma}
\end{frame}

\begin{frame}{DGL Cross-Year Message}
  \threebulletsummary
    {The main reporting target is a stable DGL result across 2022F, 2023D, and 2024I.}
    {The same tight symmetric category should be used for all three years to keep the interpretation clean.}
    {Once the remaining year outputs are filled, this section becomes the main result block of the presentation.}
\end{frame}

\section{DSA Status}

\begin{frame}{Why DSA Is Separate}
  \threebulletsummary
    {The DSA study is not yet a headline resolution result because the tag/probe bias remains visible after the symmetric downstream cuts.}
    {The bias persists even for exactly-two-DSA events and also remains in the symmetric residual diagnostic.}
    {For that reason, DSA should be presented as a status and next-step study rather than as a final resolution measurement.}
\end{frame}

\begin{frame}{DSA Status: Symmetric Residual Diagnostic}
  \twoplots
    {\DSABaseDir/compare_MC_DATA_DSA_\TightLabel/pt_mean_mc_vs_data.png}
    {DSA tight: $p_{T}$ mean, MC vs Data}
    {\DSABaseDir/compare_MC_DATA_DSA_\TightLabel/pt_sigma_mc_vs_data.png}
    {DSA tight: $p_{T}$ sigma, MC vs Data}
\end{frame}

\begin{frame}{DSA Status: Symmetric Residual Cross-Check}
  \twoplots
    {\DSABaseDir/compare_MC_DATA_DSA_\TightLabel/pt_mean_sym_mc_vs_data.png}
    {DSA tight: symmetric residual mean}
    {\DSABaseDir/compare_MC_DATA_DSA_\TightLabel/pt_sigma_sym_mc_vs_data.png}
    {DSA tight: symmetric residual sigma}
\end{frame}

\begin{frame}{DSA Findings So Far}
  \threebulletsummary
    {The probe tends to have systematically larger $p_{T}$ than the tag in DSA, and this directly explains the residual bias.}
    {The effect remains when restricting to exactly two reconstructed DSA muons, so it is not mainly a multi-candidate probe-choice artifact.}
    {The next investigation therefore moves upstream, into the ntuplizer and the standalone-track quantities themselves.}
\end{frame}

\begin{frame}{Approved DSA Next Steps}
  \begin{itemize}
    \item Keep the current tag/probe selection unchanged while debugging the upstream source of the bias.
    \item Use the dedicated ntuplizer study branch with extra DSA raw branches: \texttt{p}, \texttt{qoverp}, \texttt{qoverpError}, reference point, and side label.
    \item Study the unordered exactly-two-DSA pair summary before applying any tag/probe interpretation.
    \item Decide only after that whether the issue is mainly track parameterization, detector-side asymmetry, or something else upstream.
  \end{itemize}
\end{frame}

\section{Takeaways}

\begin{frame}{Summary}
  \begin{itemize}
    \item The immediate reporting target is a clean DGL result for 2022F, 2023D, and 2024I using one consistent selection category.
    \item The DGL control plots and MC vs Data overlays form the main physics message of the current update.
    \item DSA remains an active bias study, not yet a final resolution result.
    \item The upstream ntuplizer extension is the approved next step for isolating the DSA source.
  \end{itemize}
\end{frame}

\end{document}
